\documentclass[11pt]{article}
\usepackage[margin=1in]{geometry}
\usepackage[T1]{fontenc}
\usepackage[utf8]{inputenc}
\usepackage{lmodern}
\usepackage{enumitem}
\usepackage{xcolor}
\usepackage{hyperref}
\usepackage{titlesec}
\hypersetup{colorlinks=true, linkcolor=blue, urlcolor=blue}
\setlist[itemize]{leftmargin=1.2em}
\setlist[enumerate]{leftmargin=1.4em}
\titleformat{\section}{\Large\bfseries\color{blue!60!black}}{}{0em}{}
\titleformat{\subsection}{\large\bfseries\color{blue!50!black}}{}{0em}{}
\title{e-puck2 Beginner Guide\\\large First Steps for New Students}
\author{Prepared for Jack}
\date{July 21, 2026}

\begin{document}
\maketitle

\section*{Official starter link}
The main official GCtronic e-puck2 guide is:
\begin{itemize}
  \item \url{https://www.gctronic.com/doc/index.php?title=e-puck2}
\end{itemize}
Start there for manufacturer documentation, then use this beginner guide as the practical first-day roadmap.

\section{What the e-puck2 is}
The e-puck2 is a small differential-drive mobile robot used for robotics education and research. It is useful for teaching:
\begin{itemize}
  \item basic programming,
  \item sensing and feedback,
  \item obstacle avoidance,
  \item control and estimation,
  \item multi-robot experiments.
\end{itemize}

For a beginner, the most important mental model is simple: the robot moves because its \textbf{left and right wheels} move at chosen speeds. If both wheels move equally, it goes roughly straight. If one wheel moves faster, it turns.

\section{What beginners should learn first}
Before writing advanced code, every student should be able to explain:
\begin{enumerate}
  \item where the robot's forward direction is,
  \item which wheels create turning,
  \item what the proximity sensors are for,
  \item why simulation should come before real-hardware testing,
  \item how to stop a bad experiment quickly.
\end{enumerate}

\section{Safety rules for first-time users}
Use these rules every time:
\begin{itemize}
  \item Start in simulation whenever possible.
  \item Keep the robot on an open surface with space around it.
  \item Use low speeds for first tests.
  \item Test one change at a time.
  \item Be ready to lift the robot or switch it off if motion is wrong.
  \item Never run a new controller at full speed first.
  \item Do not assume sensor readings are perfect.
\end{itemize}

\section{Recommended learning order}
A beginner-friendly order is:
\begin{enumerate}
  \item Understand the robot body, wheels, and basic sensors.
  \item Run a simple e-puck2 simulation.
  \item Try open-loop motion: forward, stop, turn.
  \item Read sensor values.
  \item Build a simple obstacle-avoidance behavior.
  \item Only then move to the real robot.
\end{enumerate}

\section{Which simulation software to use}
For beginners, use \textbf{Webots} as the main simulation software. It is the most natural first choice for e-puck2 learning because it already supports mobile-robot simulation well and is widely used for educational robotics workflows.

Recommended links:
\begin{itemize}
  \item Webots getting started guide: \href{https://cyberbotics.com/doc/guide/getting-started-with-webots?version=R2025a}{Cyberbotics Webots guide}
  \item GCtronic e-puck2 page: \href{https://www.gctronic.com/doc/index.php?title=e-puck2}{Official e-puck2 wiki}
\end{itemize}

A simple beginner workflow is:
\begin{enumerate}
  \item Install and open Webots.
  \item Load an e-puck2 world or example project.
  \item Run the simulation and watch how the robot moves.
  \item Change a small controller constant.
  \item Re-run and confirm the behavior changed.
  \item Only after the simulation behaves well, try the same idea on the real robot at lower speed.
\end{enumerate}

\section{Basic starter code}
Students should not start with a large project. Start with a tiny controller that moves forward, turns, and stops. Even pseudocode-level understanding is useful on day one.

Example beginner Python-style controller logic:
\begin{verbatim}
# basic idea: move forward for a while, then turn, then stop

left_speed = 2.0
right_speed = 2.0

# drive forward
set_wheels(left_speed, right_speed)
wait(2.0)

# turn in place
set_wheels(-1.5, 1.5)
wait(1.0)

# stop
set_wheels(0.0, 0.0)
\end{verbatim}

What this teaches:
\begin{itemize}
  \item equal wheel speeds produce roughly straight motion,
  \item opposite wheel directions produce an in-place turn,
  \item stopping is just setting both wheels to zero.
\end{itemize}

A first reactive behavior can be explained like this:
\begin{verbatim}
front_sensor = read_front_proximity()

if front_sensor > threshold:
    set_wheels(-1.0, 1.0)   # turn away
else:
    set_wheels(2.0, 2.0)    # move forward
\end{verbatim}

This is enough to introduce the idea of \textbf{sense -> decide -> act}, which is the core robotics loop beginners need to understand.

\section{Connection modes beginners should know}
The e-puck2 can be used in several connection modes. For beginners, use \textbf{USB first}. It removes unnecessary wireless debugging problems.

\subsection{USB: recommended first path}
\begin{itemize}
  \item Selector position: \textbf{8}
  \item Why use it first: simplest setup and easiest debugging
  \item Common Linux device names: \texttt{/dev/ttyACM0}, \texttt{/dev/ttyACM1}, \texttt{/dev/ttyACM2}
\end{itemize}

Typical first-use sequence:
\begin{enumerate}
  \item Set the selector to \textbf{8}.
  \item Connect one micro-USB cable.
  \item Turn the robot on.
  \item On Linux, make sure the user has serial permission if needed:
\begin{verbatim}
sudo adduser $USER dialout
\end{verbatim}
  \item Log out and back in if the group was changed.
  \item Check whether serial devices appeared:
\begin{verbatim}
ls /dev/ttyACM*
\end{verbatim}
\end{enumerate}

\subsection{Bluetooth}
\begin{itemize}
  \item Selector position: \textbf{3}
  \item Usually not the best first choice for beginners on Linux
  \item Can work well later, but pairing and serial binding add extra steps
\end{itemize}

\subsection{WiFi}
\begin{itemize}
  \item Selector position: \textbf{15 (F)}
  \item Good when untethered control is needed
  \item Requires the correct WiFi firmware and network configuration
\end{itemize}

If the robot has no saved WiFi settings, it can create its own access point. A common first setup flow is:
\begin{enumerate}
  \item Set selector to \textbf{15 (F)}.
  \item Connect to the robot SSID \texttt{e-puck2\_0XXXX}.
  \item Use password \texttt{e-puck2robot}.
  \item Open \url{http://192.168.1.1}.
  \item Save the desired WiFi network credentials.
  \item Power-cycle the robot once.
\end{enumerate}

\section{What to test on day one}
A good first-day progression is:
\begin{enumerate}
  \item Connect to the robot over USB.
  \item Confirm the host computer can see the robot.
  \item Read sensor values without moving the robot.
  \item Blink an LED or perform another harmless check.
  \item Command a tiny forward motion.
  \item Command a tiny turn.
  \item Stop and write down what happened.
\end{enumerate}

\section{Simple beginner experiments}
These are appropriate first experiments:
\begin{itemize}
  \item Drive straight for a short time.
  \item Turn left and right in place slowly.
  \item Stop when a front obstacle is detected.
  \item Compare left and right proximity readings near a wall.
  \item Run the same controller in simulation and then on the real robot at lower speed.
\end{itemize}

\section{Common beginner mistakes}
Watch for these:
\begin{itemize}
  \item Going to real hardware too early.
  \item Testing at high speed first.
  \item Making many code changes before the next run.
  \item Forgetting to check selector position.
  \item Using a charge-only USB cable.
  \item Assuming USB port numbering will always stay the same.
  \item Confusing Bluetooth mode with WiFi mode.
\end{itemize}

\section{What students should record in their lab notes}
For each run, write down:
\begin{itemize}
  \item date and robot ID,
  \item connection mode used,
  \item code version or script name,
  \item speed settings,
  \item expected behavior,
  \item actual behavior,
  \item next change to try.
\end{itemize}

\section{A practical beginner checklist}
Before saying ``the robot is working,'' verify these basics:
\begin{itemize}
  \item the selector is in the intended position,
  \item the robot powers on correctly,
  \item the host computer detects the device or network endpoint,
  \item a sensor read succeeds,
  \item a tiny motion command works,
  \item the robot can be stopped safely.
\end{itemize}

\section{Bottom line}
For beginners, the right workflow is:\
\textbf{learn the robot body \textrightarrow{} start in simulation \textrightarrow{} connect over USB \textrightarrow{} verify sensing \textrightarrow{} try tiny motions \textrightarrow{} only then attempt more advanced behaviors.}

That sequence prevents most frustrating first-day failures and gives students a clean path from curiosity to safe real-robot control.

\end{document}
